\documentclass[11pt,a4paper]{article}

\usepackage[utf8]{inputenc}

\usepackage[T1]{fontenc}

\usepackage[polish]{babel}

\usepackage{lmodern}

\usepackage{geometry}

\geometry{a4paper,margin=27mm}

\usepackage{microtype}

\usepackage[hidelinks]{hyperref}

% Portfolio: MSC-AR-P17; source: prospective ARC-XR research scope.

\title{Kalibracja niepewności rozpoznania stanu użytkownika a selektywna abstencja adaptatora XR}

\author{Mateusz Pomianek}

\date{Wersja 0.3.0, 21 listopada 2026}

\begin{document}

\maketitle

\begin{abstract}

Przedmiotem opracowania jest Kalibracja niepewności rozpoznania stanu użytkownika a selektywna abstencja adaptatora XR. Czy abstencja przy wysokiej niepewności modelu zmniejsza koszt błędnych interwencji XR? Porównane zostaną dwie polityki adaptatora XR: obowiązkowe podjęcie decyzji oraz powstrzymanie zmiany interfejsu przy niepewnych danych. Głównym punktem końcowym jest krzywa ryzyka błędnej adaptacji względem pokrycia decyzji. Opis ma charakter wstępny i nie przedstawia niezweryfikowanych wyników.

\end{abstract}

\section{Wprowadzenie}

Rozwój adaptacyjnych systemów XR stwarza możliwość badania interakcji człowieka z informacją przestrzenną, ale samo zastosowanie AR, VR lub modeli sztucznej inteligencji nie przesądza o przewadze poznawczej czy ergonomicznej. W rozpatrywanym zagadnieniu, Kalibracja niepewności rozpoznania stanu użytkownika a selektywna abstencja adaptatora XR, kluczowe jest rozdzielenie właściwości systemu od zmienności zachowania użytkownika. Porównanie należy zaprojektować tak, aby mierzyć jednoznaczne zjawisko, a nie tylko ogólną ocenę technologii.

\section{Cel pracy i pytanie badawcze}

Czy abstencja przy wysokiej niepewności modelu zmniejsza koszt błędnych interwencji XR? Projekt wymaga krytycznego przeglądu aktualnego stanu badań oraz zdefiniowania nowości w stosunku do innych tematów portfolio XR.

\section{Zakres i odrębny wkład naukowy}
Selektywna kalibracja modelu rozpoznającego stan użytkownika i abstencja adaptatora przy niepewności. Badanie dotyczy jakości decyzji modelu, nie behawioralnego kosztu odmowy w MSC-AR-P25. Walidacja niepewności i selekcji, nie behawioralny wynik użytkownika w programie P25.
\section{Hipotezy badawcze}
\paragraph{H1.} Kalibracja niepewności wraz z abstencją zmniejsza ryzyko błędnej adaptacji przy ustalonym pokryciu decyzji.
\paragraph{H2.} Zysk utrzymuje się na niezależnych sesjach, bez dobierania progu i kalibratora po analizie końcowej.
\paragraph{H3.} Metoda kalibracyjna utrzymuje ograniczenia błędu w podgrupach nowych użytkowników, a nie jedynie w próbie średniej.
\section{Projekt badania i plan eksperymentów}
\subsection{A. Operacjonalizacja i kwalifikacja danych}
Selektywna kalibracja modelu rozpoznającego stan użytkownika i abstencja adaptatora przy niepewności. Badanie dotyczy jakości decyzji modelu, nie behawioralnego kosztu odmowy w MSC-AR-P25. Jednostka wnioskowania: sesja testowa i błędna decyzja adaptacyjna w jej obrębie. W protokole zamrozić definicję i mianownik wyniku oraz źródło prawdy odniesienia, zgodność znaczników czasu, prawa do danych i zasady rekrutacji.
\subsection{B. Eksperyment porównawczy}
Na identycznym zamrożonym modelu porównać wymuszony wybór, próg globalny i selekcję zależną od kalibracji; dane kalibracyjne rozdzielić od testu. Estymand pierwszorzędowy: pole pod krzywą ryzyko–pokrycie decyzji dla kalibrowanego adaptatora vs adaptatora bez selekcji. Przed testem zdefiniować randomizację, próbę, punkt końcowy i warunki wykluczenia bez selekcji wyniku post hoc.
\subsection{C. Kontrole i analiza odporności}
Definiować błędną adaptację względem niezależnej oceny zadania, a nie przekonania modelu o poprawności własnej decyzji. Raportować ryzyko wraz z pokryciem i odsetkiem abstencji, aby algorytm nie wygrywał przez wstrzymywanie wszystkiego. W rejestrze uwzględnić braki, wstrzymania, opóźnienia i niepełne próby zamiast usuwać je z analizy.
\subsection{D. Walidacja na niezależnym materiale}
Weryfikować na niezależnych sesjach lub osobach, bez dostępu modelu do danych końcowych. Przedziały krzywej ryzyko–pokrycie bootstrapem po sesjach; konfrontacja na stałym minimalnym pokryciu. Dokumentować wyniki niesprzyjające hipotezie oraz rzeczywiste granice przeniesienia.
\section{Estymand, analiza statystyczna i decyzje}
pole pod krzywą ryzyko–pokrycie decyzji dla kalibrowanego adaptatora vs adaptatora bez selekcji. Jednostka: sesja testowa i błędna decyzja adaptacyjna w jej obrębie. Przedziały krzywej ryzyko–pokrycie bootstrapem po sesjach; konfrontacja na stałym minimalnym pokryciu. H1 jest kandydatem hipotezy potwierdzającej, a H2/H3 pozostają eksploracyjne do formalnego planu wielokrotnych porównań. Liczebność uczestników, margines, próg i moc lub precyzja wymagają pilotażu oraz zamrożenia przed testem końcowym.
\section{Stan badań i granica nowości}
\begin{itemize}
\item Existing adaptive physiological feedback in AR means adapting an interface is not by itself a novel contribution. \cite{xrprior1}.
\item Physiology-driven VR adaptation has been demonstrated; independent behavioural gains require direct controlled trials. \cite{xrprior2}.
\end{itemize}
Walidacja niepewności i selekcji, nie behawioralny wynik użytkownika w programie P25. Odrębność od wcześniejszych publikacji i innych tematów programu należy sprawdzić przed rejestracją badania.
\section{Warunki bezpieczeństwa, sprzętu i zgodności}
Jest to plan prospektywny, nie dowód wykonania eksperymentu lub zatwierdzonej etyki. Udział osób, rejestrowanie biometrii i interwencje wymagają uprzedniej procedury etycznej i świadomej zgody. Nie wnioskować o emocjach ani zdrowiu z samych sygnałów. Standardowy Quest 3 nie ma okulometrii HMD, a Tobii Eye Tracker 5 służy do stanowiska ekranowego; każde połączenie zweryfikować przed protokołem.
\section{Materiał i metody}

Porównane zostaną dwie polityki adaptatora XR: obowiązkowe podjęcie decyzji oraz powstrzymanie zmiany interfejsu przy niepewnych danych. Wymagane jest zdefiniowanie jednostki analizy, zasad doboru przypadków, wersji scenariusza i źródeł pomiaru. Rozkład warunków oraz reguły wykluczenia prób należy określić przed pomiarem. W badaniach z udziałem ludzi niezbędna jest uprzednia właściwa ocena etyczna i świadoma zgoda. W testach wyłącznie symulacyjnych należy oddzielić prawdę symulacyjną od obserwacji empirycznej.

\subsection{Wynik główny i kryteria jakości}

Krzywa ryzyka błędnej adaptacji względem pokrycia decyzji. Pomiar powinien zachowywać jawne jednostki, znaczniki czasu i reguły oceny poprawności. Przy modelach predykcyjnych stosować rozłączne partycje oceny; przy eksperymentach interwencyjnych oddzielić skuteczność predykcji od przyczynowego wpływu interfejsu.

\section{Plan analizy}

Ocena selective-risk i coverage na rozłącznych scenariuszach testowych i uczestnikach; raportować także liczbę decyzji wstrzymanych. W wynikach należy raportować wielkość efektu lub błąd estymacji z przedziałem niepewności, liczbą obserwacji i odsetkiem odrzuconych prób. Analizy dodatkowe wymagają jasnego oznaczenia jako eksploracyjne. Nie wybierać progów decyzyjnych ani modelu na podstawie obserwacji ze zbioru testowego.

\section{Wyniki}

W obecnej wersji nie raportuje się wyników pomiarów, tabel statystycznych ani parametrów wyznaczonych z nieudokumentowanych danych. To kierunek badawczy, a nie zatwierdzony protokół badania ani zgłoszony artykuł.

\section{Dyskusja}

Interpretacja przyszłych wyników musi uwzględnić alternatywne mechanizmy i zależność wniosków od jakości danych, warunków zadania oraz ograniczeń sprzętu. Wartość potencjalnego rozwiązania zależy nie tylko od trafności modelu, lecz także od jego opóźnienia, stabilności i wpływu na zachowanie człowieka. Wyniki testów syntetycznych nie mogą być automatycznie uogólniane na populację uczestników.

\section{Ograniczenia badania}

Reguły progowe muszą zostać określone bez wykorzystania zbioru testowego; niskie ryzyko przy znikomym pokryciu nie wystarcza. Metoda selekcji / ocena prawdopodobieństwa abstencji; nie efekt behawioralny odrzucenia adaptacji. Żadna część szkicu nie dowodzi przeprowadzenia zbierania danych, uzyskania zgód etycznych ani gotowości manuskryptu do zgłoszenia.

\section{Wnioski}

Przedstawiono problem, kontrast i sposób jego planowanej oceny. Rozstrzygnięcie postawionego pytania wymaga przygotowania protokołu, zweryfikowanych źródeł oraz poprawnie udokumentowanych wyników. Bez tych materiałów nie można sformułować konkluzji empirycznych.

\begin{thebibliography}{9}
\bibitem{xrprior1} Designing adaptive feedback systems for managing cognitive load in augmented reality (2025), Design Science 11, e49. \url{https://doi.org/10.1017/dsj.2025.10040}.
\bibitem{xrprior2} Towards Intelligent VR Training: A Physiological Adaptation Framework for Cognitive Load and Stress Detection (2025), UMAP 2025. \url{https://doi.org/10.1145/3699682.3727575}.
\end{thebibliography}
\end{document}

